% !TEX program = xelatex
% 공통수학2 · Ⅱ. 집합과 명제 · 2. 명제 · 03 명제의 증명 · 2/2차시
\documentclass[11pt,a4paper]{article}
\usepackage{kotex}
\usepackage{iftex}
\ifPDFTeX\else
  \setmainhangulfont{Noto Serif CJK KR}
  \setsanshangulfont{Noto Sans CJK KR}
\fi
\usepackage[margin=2.1cm]{geometry}
\usepackage{parskip}
\usepackage{enumitem}
\usepackage{array}
\usepackage{tabularx}
\usepackage{booktabs}
\usepackage{xcolor}
\usepackage{amssymb}
\usepackage{tikz}
\usepackage[most]{tcolorbox}
\usepackage{fancyhdr}
\usepackage{titlesec}

\definecolor{goldc}{HTML}{B07C22}
\definecolor{goldstrong}{HTML}{8A5F16}
\definecolor{indigoc}{HTML}{2B3B57}
\definecolor{indigosoft}{HTML}{6E82A6}
\definecolor{linec}{HTML}{D6DACB}
\definecolor{surface2}{HTML}{E4E8DC}
\definecolor{notebg}{HTML}{FBF3E3}
\definecolor{inksoft}{HTML}{52604F}

\titleformat{\section}{\normalfont\sffamily\bfseries\large\color{indigoc}}{\thesection}{0.6em}{}
\titlespacing{\section}{0pt}{16pt}{8pt}
\newtcolorbox{ideabox}{enhanced, breakable, colback=white, colframe=goldc, boxrule=0.5pt, leftrule=3pt, arc=2pt, left=10pt, right=10pt, top=8pt, bottom=8pt}
\newtcolorbox{calloutbox}{enhanced, breakable, colback=white, colframe=indigosoft, boxrule=0.5pt, leftrule=3pt, arc=2pt, left=10pt, right=10pt, top=7pt, bottom=7pt}
\newtcolorbox{notebox}{enhanced, breakable, colback=notebg, colframe=goldc, boxrule=0.5pt, leftrule=3pt, arc=2pt, left=10pt, right=10pt, top=7pt, bottom=7pt}
\newtcolorbox{answerbox}{enhanced, breakable, colback=white, colframe=linec, boxrule=0.5pt, arc=2pt, left=10pt, right=10pt, top=7pt, bottom=7pt}
\newcommand{\lessonbanner}[3]{%
  \begin{tcolorbox}[enhanced, colback=indigoc, colframe=indigoc, boxrule=0pt, arc=0pt,
    left=14pt, right=14pt, top=14pt, bottom=10pt, borderline south={2.4pt}{0pt}{goldc}, fontupper=\color{white}]
    {\sffamily\footnotesize\color{white!85!black} #1}\\[4pt]{\sffamily\bfseries\Large #2}\\[3pt]{\sffamily\small\color{white!88!black} #3}
  \end{tcolorbox}}
\pagestyle{fancy}\fancyhf{}\renewcommand{\headrulewidth}{0pt}
\fancyfoot[L]{\footnotesize\color{inksoft} 공통수학2 · 2026학년도 2학기 · 지평선고등학교}
\fancyfoot[R]{\footnotesize\color{inksoft} 교사용 지도안 -- 배포 금지}

\begin{document}
\lessonbanner{공통수학2 · Ⅱ. 집합과 명제 · 2. 명제}{03 명제의 증명 --- 2/2차시}{절대부등식 · 교사용 지도안 (2026학년도 2학기)}
\vspace{10pt}

\noindent\begin{tabularx}{\linewidth}{@{}>{\bfseries\color{indigoc}}p{2.6cm}X@{}}
\toprule
대상 & 고1 · 2개 반 · 반당 13명 \\
차시 & 50분 (도입 10 · 전개 30 · 정리 10) \\
성취기준 & 절대부등식의 뜻을 알고, 간단한 절대부등식을 증명할 수 있다. \\
선행 학습 & 30차시 --- 대우를 이용한 증명법과 귀류법 \\
\bottomrule
\end{tabularx}

\section{학습 목표 \& Big Idea}
\begin{ideabox}
\textbf{\color{goldstrong}영속적 이해 (Big Idea)}\\
`모든 실수 $x$에 대하여 항상 성립하는' 절대부등식은 어떤 실수의 제곱도 0 이상이라는 단 하나의 사실($a^2\ge0$)로부터 완전제곱식을 만들어 증명할 수 있다.
\vspace{4pt}
\small\color{inksoft} 핵심개념: \textbf{절대부등식} \quad\textbullet\quad 핵심개념: \textbf{완전제곱식} \quad\textbullet\quad 일반화: \textbf{차를 완전제곱 꼴로 만들면 부등식이 증명된다}
\end{ideabox}
\begin{calloutbox}
\textbf{차시 학습 목표}
\begin{itemize}[leftmargin=1.4em, itemsep=2pt, topsep=4pt]
  \item 절대부등식의 뜻을 알고, 완전제곱식을 이용하여 간단한 절대부등식을 증명할 수 있다.
  \item 산술평균과 기하평균의 관계를 이해하고 이를 이용하여 절대부등식을 증명할 수 있다.
\end{itemize}
\end{calloutbox}

\section{질문의 연결}
\begin{tabularx}{\linewidth}{@{}>{\bfseries\color{indigoc}\small}p{2.4cm}X@{}}
사실적 질문 & 어떤 부등식을 `절대부등식'이라고 부를 수 있을까? \\[6pt]
개념적 질문 & 왜 $A-B\ge0$을 완전제곱식들의 합으로 나타낼 수 있으면 부등식 $A\ge B$가 항상 성립함이 증명될까? \\[6pt]
논쟁적 질문 & 등호가 성립하는 경우까지 정확히 찾아야만 증명이 `완전하다'고 할 수 있을까, 아니면 부등식만 보이면 충분할까? \\
\end{tabularx}

\section{수업 흐름}
\begin{tabularx}{\linewidth}{@{}>{\centering\arraybackslash}p{1.6cm}X>{\raggedright\arraybackslash\small}p{3.1cm}@{}}
\toprule
\textbf{단계} & \textbf{활동 \& 발문} & \textbf{ATL / VTR} \\
\midrule
\textcolor{goldstrong}{\textbf{도입}}\newline\footnotesize 10분 &
\textbf{마음공부 (2분)} --- 유무념 대조.\newline
\textbf{생각 열기 (8분)} --- $x^2+1>0$, $x^2-1>0$, $|x-1|\ge0$ 중 모든 실수 $x$에 대하여 항상 성립하는 것 찾기($x=0$을 대입해 반례 확인). &
ATL: 사고(반례 대입)\newline VTR: See-Think-Wonder \\
\addlinespace
\textcolor{goldstrong}{\textbf{전개}}\newline\footnotesize 30분 &
\textbf{절대부등식과 기본 성질 (8분)} --- 절대부등식 정의 $\to$ 실수의 기본 성질($a>b \Leftrightarrow a-b>0$, $a^2\ge0$, $a^2+b^2\ge0$, $a^2+b^2=0 \Leftrightarrow a=b=0$).\newline
\textbf{완전제곱식을 이용한 증명 (12분)} --- 예제 3($a^2+b^2\ge ab$, $\{a-\frac b2\}^2+\frac34b^2$ 꼴로 변형) $\to$ 문제 3($(a+b)^2\ge4ab$; $a^2-4ab+5b^2\ge0$), 등호 성립 조건 강조.\newline
\textbf{산술평균과 기하평균 (10분)} --- 예제 4($a>0,b>0$일 때 $\frac{a+b}2\ge\sqrt{ab}$, $(\frac{\sqrt a}{}-\sqrt b)^2\ge0$으로 유도) $\to$ 문제 4($a+\frac1a\ge2$). &
ATT: 촉진자 역할\newline ATL: 대수적 조작·논리 전개 \\
\addlinespace
\textcolor{goldstrong}{\textbf{정리}}\newline\footnotesize 10분 &
\textbf{개념 연결 질문 (5분)} --- ``등호가 성립하는 조건까지 찾아야 증명이 완전한 이유는 무엇일까?''\newline
\textbf{성찰 (5분)} --- 감사 일기 1문장 + `명제의 증명' 및 대단원 Ⅱ `집합과 명제' 마무리 + 다음 차시 예고(중단원·대단원 마무리). &
마음공부 · 감사 일기 \\
\bottomrule
\end{tabularx}

\begin{calloutbox}
\textbf{지도상의 유의점} --- 대소 비교를 위해 부등식의 양변을 제곱할 때는 반드시 두 변이 음이 아닌지 확인하도록 지도한다. 등호가 포함된 부등식을 증명할 때는 특별한 언급이 없어도 등호가 성립하는 조건을 찾아 증명 과정에 포함하게 한다.
\end{calloutbox}

\section{형성평가 \& 답안}
\begin{minipage}[t]{0.48\linewidth}
\begin{answerbox}
\textbf{문제 3} \quad 완전제곱식 이용\\
\small ⑴ $(a+b)^2\ge4ab$\\
⑵ $a^2-4ab+5b^2\ge0$\\[4pt]
\textbf{\color{goldstrong}답} \; ⑴ $(a+b)^2-4ab=(a-b)^2\ge0$, 등호는 $a=b$일 때. \; ⑵ $a^2-4ab+5b^2=(a-2b)^2+b^2\ge0$, 등호는 $a=2b$이고 $b=0$ 즉 $a=b=0$일 때.
\end{answerbox}
\end{minipage}\hfill
\begin{minipage}[t]{0.48\linewidth}
\begin{answerbox}
\textbf{문제 4} \quad 산술평균·기하평균 ($a>0$)\\
\small $a+\dfrac1a\ge2$\\[4pt]
\textbf{\color{goldstrong}답} \; $a>0$이고 $\frac1a>0$이므로 산술평균과 기하평균의 관계로부터 $a+\frac1a\ge2\sqrt{a\cdot\frac1a}=2$. 등호는 $a=\frac1a$, 즉 $a=1$일 때 성립.
\end{answerbox}
\end{minipage}

\section{성취수준 (이 차시 범위)}
\begin{tabularx}{\linewidth}{@{}>{\bfseries\color{goldstrong}}p{0.9cm}X@{}}
\toprule
A & 절대부등식의 뜻을 설명하고, 완전제곱식이나 산술·기하평균 관계를 이용하여 간단한 절대부등식을 스스로 증명할 수 있다. \\
B & 절대부등식의 뜻을 알고, 안내된 절차를 따라 절대부등식을 증명할 수 있다. \\
C & 절대부등식의 뜻을 알고, 간단한 절대부등식이 성립함을 확인할 수 있다. \\
D & 절대부등식의 뜻을 안다. \\
E & 부등식과 절대부등식을 구별할 수 있다. \\
\bottomrule
\end{tabularx}

\section{다음 차시}
\begin{calloutbox}
\textbf{중단원·대단원 마무리 --- 2. 명제 / Ⅱ. 집합과 명제.} `명제와 조건', `명제 사이의 관계', `명제의 증명'을 통합 정리하고, 대단원 Ⅱ `집합과 명제' 전체를 마무리한다.
\end{calloutbox}
\end{document}
